\section{Proofs for Section~\ref{sec:accounting}}\label{app:proofs-accounting}

Every statement below is checked numerically in \texttt{code/verify\_part2\_accounting.py}; check numbers are given in brackets.

\subsection*{Proof of Proposition~\ref{prop:acct-path}}

(a) Under Assumption~\ref{ass:acct-books}(iv) the cash score records the cash flows of Assumption~\ref{ass:tasks-transition} on the success path: the flow $bW\varrho(t)-C^{\mathrm{par}}$ on $[0,T_{\mathrm{par}})$ and $bW$ afterwards, $F+I+C^J$ at date 0 and $C^{\mathrm{str}}=C^{\mathrm{sev}}+C^{\mathrm{brk}}$ at $T_{\mathrm{par}}$. None of $F^{\mathrm{cap}}$, $T_B$, $t_{\mathrm{go}}$, $t_{\mathrm{sev}}$, $t_{\mathrm{com}}$, $t_{\mathrm{brk}}$ or the classification of outflows enters, and integrating the flow of Proposition~\ref{prop:tasks-scorepath}(a) over $[y-1,y)$ and adding the one-time charges that fall in that year gives $\Delta R^c_y$. [1a]

(b) Compare the two bases item by item, cumulatively to the end of year $Y\ge 1$, and multiply by $S$. The operating flow and $C^J$ are charged identically. The build is charged $F$ at date 0 in cash, and $F^{\mathrm{exp}}$ at date 0 plus $F^{\mathrm{cap}}\mathsf A_Y$ of amortization in GAAP, so GAAP charges $F^{\mathrm{cap}}(1-\mathsf A_Y)$ less. The equipment is charged $I$ in cash and nothing beyond the depreciation already netted in $W$ in GAAP, by (ii). The legacy assets are charged $C^{\mathrm{wo}}\min\{Y/T_{\mathrm{par}},1\}$ in GAAP and nothing in cash. Restructuring is charged when paid in cash and when recognized in GAAP, so GAAP charges more by the recognized but unpaid amount $\mathsf L_Y$. Since a lower charge raises the score, summing gives \eqref{eq:acct-accrual}. When $\mathsf A_Y=1$, $\mathsf L_Y=0$ and $Y\ge T_{\mathrm{par}}$, the right side is $I-C^{\mathrm{wo}}$. [1b]

(c) If $T_{\mathrm{par}}\le 1\le T_B$, the first year contains the date-0 charges, the whole parallel run, the cutover and, by assumption, all restructuring charges in both bases, so restructuring creates no difference. Amortization in year 1 is $F^{\mathrm{cap}}|[0,1)\cap[t_{\mathrm{go}},t_{\mathrm{go}}+T_B)|/T_B=F^{\mathrm{cap}}(1-t_{\mathrm{go}})/T_B$, because $t_{\mathrm{go}}+T_B\ge 1$. The year-1 difference is therefore $F-F^{\mathrm{exp}}-F^{\mathrm{cap}}(1-t_{\mathrm{go}})/T_B+I-C^{\mathrm{wo}}=F^{\mathrm{cap}}[1-(1-t_{\mathrm{go}})/T_B]+I-C^{\mathrm{wo}}$, divided by $S$. With $F^{\mathrm{cap}}=F$, $t_{\mathrm{go}}=0$ and $T_B=T_A$ it is $(F+I-F/T_A-C^{\mathrm{wo}})/S$, the difference in Proposition~\ref{prop:tasks-scorepath}(c). The coefficient $1-(1-t_{\mathrm{go}})/T_B$ lies in $[0,1)$ and increases with $t_{\mathrm{go}}$ and $T_B$. If $C^{\mathrm{sev}}$ is recognized in year 0, the GAAP charge moves from year 1 to year 0, and by (a) the cash path is unchanged. [1c, 1e, 1h]

(d) For $y\ge 2$ no date-0 charge, cutover or restructuring charge falls in year $y$, $C^{\mathrm{wo}}$ has been charged by $T_{\mathrm{par}}\le 1$, and the cash flow is $bW$ throughout the year; the only difference is the amortization in year $y$. It vanishes after $t_{\mathrm{go}}+T_B\ge 1\ge T_{\mathrm{par}}$. [1d]

(e) Costs incurred while significant development uncertainty exists do not meet the probable-to-complete threshold (ASC 350-40-25-12A), and data conversion and training costs are expensed (ASC 350-40-25-4 and 25-5), so at most $(1-u)(F-F^{\mathrm{dc}})$ can be capitalized. The gap in \eqref{eq:acct-first} is nondecreasing in $F^{\mathrm{cap}}$, which gives the bound; at $u=1$ the bound is $(I-C^{\mathrm{wo}})/S$. [1f] \qed

\subsection*{Proof of Proposition~\ref{prop:acct-program}}

(a) The vintage spent at $s$ is amortized at the rate $C^{\mathrm{cap}}(s)/T_B$ on $[s+L,s+L+T_B)$. At $t$ the charge is
\[
\frac1{T_B}\int_{t-L-T_B}^{t-L}C^{\mathrm{cap}}_0e^{g_Fs}\,ds=C^{\mathrm{cap}}(t)\,e^{-g_FL}\,\frac{1-e^{-g_FT_B}}{g_FT_B}=C^{\mathrm{cap}}(t)\,e^{-g_FL}\mathsf A(g_FT_B),
\]
and the cash basis charges $C^{\mathrm{cap}}(t)$; dividing the difference by $S(t)$ gives \eqref{eq:acct-wedge}. [2a]

(b) Write $\phi(g):=e^{-gL}\mathsf A(gT_B)=\int_0^1e^{-g(L+T_Bv)}\,dv$, using $\mathsf A(x)=\int_0^1e^{-xv}dv$. Then $\phi(0)=1$ and $\phi'(g)=-\int_0^1(L+T_Bv)e^{-g(L+T_Bv)}dv<0$, so $1-\phi(g)$ has the sign of $g$. For $g>0$, $\partial\phi/\partial T_B=-g\int_0^1ve^{-g(L+T_Bv)}dv<0$ and $\partial\phi/\partial L=-g\phi<0$, so $w$ increases with $T_B$ and $L$; $\phi>0$ gives $w<C^{\mathrm{cap}}/S$; and $\phi(g)=1-g\int_0^1(L+T_Bv)dv+O(g^2(L+T_B)^2)=1-g(L+T_B/2)+O(g^2(L+T_B)^2)$. [2b]

(c) With $C^{\mathrm{cap}}(t)=c^{\mathrm{cap}}S(t)$ and $g_F=g$, $w$ is the constant $c^{\mathrm{cap}}[1-e^{-gL}\mathsf A(gT_B)]$; growth is the same on both bases, so the score difference equals the margin difference. [2d] \qed

\subsection*{Proof of Proposition~\ref{prop:acct-tax}}

(a) Deductions are $1/(2T)$ in tax year 1, $1/T$ in years 2 to $T$ and $1/(2T)$ in year $T+1$, and the deduction of year $y$ is discounted by $(1+r)^{-(y-1)}$. Summing the geometric series gives the closed form. The deduction of year $y$ is $P(U\in[y-1,y))$ for $U$ uniform on $[1/2,T+1/2)$, so $\psi_T=\mathbb E[(1+r)^{-\lfloor U\rfloor}]$. Raising $T$ replaces $U$ by a variable that strictly first-order stochastically dominates it, and $(1+r)^{-k}$ is strictly decreasing in $k$, so $\psi_T$ falls; every term with $k\ge1$ falls with $r$. At $r=0.10$, $\psi_5=0.796$ and $\psi_{15}=0.532$. [3a, 3b]

(b) Let $O$ be the first-year pre-tax cash change before the build and the equipment, so that $S\Delta R^c_1=O-F-I$, and every component of $O$ is deductible when paid. The first-year deduction for the build and equipment is $(F-F^{\mathrm{for}})+F^{\mathrm{for}}/30+I$: domestic expenditure is deducted when paid (\S174A(a)), foreign expenditure is amortized over 15 years from the midpoint of the year (\S174(a)), so $1/30$ falls in the first year, and the equipment takes the full allowance. With full loss offset, taxes change by $\tau[O-(F-F^{\mathrm{for}})-F^{\mathrm{for}}/30-I]=\tau[S\Delta R^c_1+F^{\mathrm{for}}(1-1/30)]$, and subtracting this from $S\Delta R^c_1$ gives the result. [3c, 3d]

(c) A dollar of domestic build is deducted in the first year, worth $\tau$; a dollar of foreign build is deducted on the 15-year schedule, worth $\tau\psi_{15}$. [3b]

(d) Capitalization and amortization in the books change neither cash flows nor taxable income. The difference between the book and tax bases is a temporary difference, recorded as deferred tax under ASC 740, which is not a cash flow. Part I's $\psi$ is the present value of deductions per dollar of capital, which is 1 for amounts deducted when incurred and $\psi_{15}$ on the foreign schedule. \qed

\subsection*{Proof of Proposition~\ref{prop:acct-cutover}}

(a) The open set $[T_{\mathrm{cl}},1-T_{\mathrm{op}}]$ has length $1-T_{\mathrm{op}}-T_{\mathrm{cl}}$. A process ready at $z<T_{\mathrm{cl}}$ waits $T_{\mathrm{cl}}-z$, one ready at an open date does not wait, and one ready at $z>1-T_{\mathrm{op}}$ waits $1+T_{\mathrm{cl}}-z$. With $z$ uniform,
\[
\mathbb E[\text{delay}]=\int_0^{T_{\mathrm{cl}}}(T_{\mathrm{cl}}-z)\,dz+\int_{1-T_{\mathrm{op}}}^1(1+T_{\mathrm{cl}}-z)\,dz=\frac{T_{\mathrm{cl}}^2}2+T_{\mathrm{op}}T_{\mathrm{cl}}+\frac{T_{\mathrm{op}}^2}2=\frac{(T_{\mathrm{op}}+T_{\mathrm{cl}})^2}2 .
\]
[4a]

(b) A material weakness is unremediated at the year-end if and only if a deficiency is introduced, it is a material weakness, and the detection time $D$ satisfies $D+T_{\mathrm{rem}}+T_{\mathrm{op}}>H$. By independence and $D\sim\mathrm{Exp}$ with mean $T_{\mathrm{dc}}$, the probability is $p_{\mathrm{def}}p_{\mathrm{mw}}\Pr(D>H-T_{\mathrm{rem}}-T_{\mathrm{op}})$, which is \eqref{eq:acct-pmw}. It is nonincreasing in $H=1-z$, and on the open set $H$ is largest at $z=T_{\mathrm{cl}}$. [4b]

(c) Cutting over at $z$ costs $L_{\mathrm{mw}}P_{\mathrm{mw}}(1-z)$ in expectation. Waiting for the next close costs $\hat c(1+T_{\mathrm{cl}}-z)$ and $L_{\mathrm{mw}}P_{\mathrm{mw}}(1-T_{\mathrm{cl}})$. Waiting is better if and only if \eqref{eq:acct-defer} holds. Write $G(z)$ for its left side minus its right side. $P_{\mathrm{mw}}(1-z)$ is nondecreasing in $z$ by (b), and $-\hat c(1+T_{\mathrm{cl}}-z)$ is strictly increasing for $\hat c>0$, so $G$ is strictly increasing and changes sign at most once. Set $z^*:=\inf\{z\in[T_{\mathrm{cl}},1-T_{\mathrm{op}}]:G(z)>0\}$, or $z^*=1-T_{\mathrm{op}}$ if the set is empty. [4c, 4d] \qed

\subsection*{Proof of Corollary~\ref{cor:acct-npv}}

Take the three branches in the proof of Proposition~\ref{prop:tasks-transition}(a). Every branch pays $C^{\mathrm{ctl}}$ at date 0. On success (probability $\varpi_p$) the control cost changes by $\Delta C^{\mathrm{ctl}}$ a year from $T_{\mathrm{par}}$ to $T_h$, worth $\Delta C^{\mathrm{ctl}}(e^{-rT_{\mathrm{par}}}-e^{-rT_h})/r$. The success branch and the undetected-failure branch (probability $(1-\varpi_p)p_u$) both cut over at $T_{\mathrm{par}}$ and bear the expected material-weakness cost $L_{\mathrm{mw}}P_{\mathrm{mw}}(H)$, valued at cutover; a detected failure is abandoned before cutover and bears none. Taking expectations and adding to \eqref{eq:tasks-npv} gives \eqref{eq:acct-npv}. [4e] \qed

\subsection*{Proof of Proposition~\ref{prop:acct-amin}}

(a) Without the shared cause, all $n_x$ layers of type $x$ fail if and only if the common-mode event of that type occurs, or it does not and each layer fails independently; the probability is $p^{\mathrm{cm}}_x+(1-p^{\mathrm{cm}}_x)p_x^{n_x}$. Types are independent given no shared cause, so the probability that every layer fails is the product, and the shared cause, independent of the rest, defeats every layer with probability $p^{\mathrm{cm}}_0$; hence $P=p^{\mathrm{cm}}_0+(1-p^{\mathrm{cm}}_0)\prod_x[\cdot]$. For one type, $P_n-P_{n+1}=(1-p^{\mathrm{cm}}_0)(1-p^{\mathrm{cm}})p^n(1-p)>0$ and $P_n\to p^{\mathrm{cm}}_0+(1-p^{\mathrm{cm}}_0)p^{\mathrm{cm}}$ for $p<1$. Each factor of a used type is at least $p^{\mathrm{cm}}_x$, and $P$ is increasing in the product. [5a, 5b]

(b) Write $Q$ for the product over the other types. The new design has $P'=p^{\mathrm{cm}}_0+(1-p^{\mathrm{cm}}_0)p^{\mathrm{ind}}P_{n-1}Q$ and the old $P=p^{\mathrm{cm}}_0+(1-p^{\mathrm{cm}}_0)P_nQ$, so $P'<P$ if and only if $p^{\mathrm{ind}}<P_n/P_{n-1}$ whatever $p^{\mathrm{cm}}_0<1$ and $Q>0$. Since $p\le 1$, $P_n=p^{\mathrm{cm}}+(1-p^{\mathrm{cm}})p^n\ge p\,p^{\mathrm{cm}}+(1-p^{\mathrm{cm}})p^n=pP_{n-1}$, so $P_n/P_{n-1}\ge p>p^{\mathrm{ind}}$ suffices. If $p^{\mathrm{cm}}>0$, $P_n$ and $P_{n-1}$ both tend to $p^{\mathrm{cm}}$, so the ratio tends to 1. [5c, 5d]

(c) $\Pi_o=1-(1-P_o)^{N_o}$ is strictly increasing in $P_o$, so $\Pi_o\le\bar\epsilon$ if and only if $P_o\le1-(1-\bar\epsilon)^{1/N_o}=\bar P_o$. The floor is then a minimum of a cost over a set, and by Definition~\ref{def:acct-admissible} the set excludes every design with $\Pi_o>\bar\epsilon$ for some $o$, whatever its cost; no loss term enters, so no saving in control cost can admit an excluded design. By (a), a design's $P_o$ is at least the common-mode floor of the types it uses at every depth, so a design whose floor exceeds $\bar P_o$ is never admissible; for an all-AI design on one provider the floor is $p^{\mathrm{cm}}_0+(1-p^{\mathrm{cm}}_0)p^{\mathrm{cm}}$, and $\bar P_o\approx\bar\epsilon/N_o$ for small $\bar\epsilon$. [5e, 5f, 5i]

(d) For one type, $P\le\bar P_o$ if and only if $p^n\le\rho$. If $\rho\le0$, no $n$ satisfies it. Otherwise $\rho<1$ because $\bar P_o<1$, and since $\ln p<0$ the condition is $n\ge\ln\rho/\ln p$, whose smallest integer solution is $\lceil\ln\rho/\ln p\rceil\ge1$. $\bar P_o$ falls with $N_o$ and rises with $\bar\epsilon$, so $\rho$ does too, and $n_{\min}$ is nondecreasing in $N_o$ and nonincreasing in $\bar\epsilon$. Depth enters only through $p^n$, which multiplies $1-p^{\mathrm{cm}}$; the common-mode terms are lowered only by another type, by (b), or by ITGC quality, by (e). [5e]

(e) $P$ is nondecreasing in each $p^{\mathrm{cm}}_x$ by (a), so a lower $p^{\mathrm{cm}}_x(\kappa_{\rm itgc})$ weakly lowers $P_o(d)$ for every design $d$ and the admissible set at a higher quality contains that at a lower one. At a given cost of the ITGCs the floor is then a minimum over a larger set, so it does not rise; minimizing jointly over designs and $\kappa_{\rm itgc}$, with the cost of the ITGCs in the objective, gives at most the floor at any fixed quality. If $p^{\mathrm{cm}}_x$ is held at its value for weak ITGCs while the cost of effective ITGCs is charged, the admissible set is too small and the cost too high. [5g, 5j]

(f) $\bar P_o$ and $P_o(d)$ depend on $N_o$, $\bar\epsilon$ and the layer parameters, not on $S$, so with $N_o$ fixed the admissible set and the minimizing cost $K^*$ do not depend on $S$, and $a_{\min}=K^*/S$, positive since $C^{\mathrm{acct}}>0$. If $N_o=oS$, $\bar P_o\to0$ as $S\to\infty$, while every design has $P_o$ at least its positive common-mode floor, so each fails at some scale. [5h] \qed

\subsection*{Proof of Proposition~\ref{prop:acct-itgc}}

(a) By Definition~\ref{def:acct-scope}, $\mathcal U^*=\mathcal U_T\cup\mathcal U_S\cup\mathcal U_I\cup\mathcal U_P$, the sets of units meeting each channel. Enlarging $\mathcal Y^*$ enlarges $\mathcal U_S$ and $\mathcal U_I$; adding an integration that feeds $\mathcal Y^*$ enlarges $\mathcal U_I$; adding ICFR-relevant policies enlarges $\mathcal U_P$; and $\mathcal U_T$ does not shrink in any case, so the union grows. A unit outside $\mathcal U_T$ is in the union if and only if it is in $\mathcal U_S\cup\mathcal U_I\cup\mathcal U_P$. An agent $g$ that performs $u$, taking over its inputs, outputs, integrations and governing policies, processes the same transactions or performs the same control (T), reads and writes at least the systems $u$ read and wrote (S), operates the integrations $u$ created or changed (I) and is governed by the same policies (P). Its channel set therefore contains that of $u$, which is nonempty when $u\in\mathcal U^*$, so $g\in\mathcal U^*$; being an AI application, $g\in\mathcal U^*_{\mathrm{AI}}$. Access the agent adds can only add channels. [6a, 6b]

(b) In Assumption~\ref{ass:acct-layers}, the term $\mathbf 1\{\text{AI used}\}(C^{\mathrm{itgc}}+C^{\mathrm{val}}/T_{\mathrm{up}})$ is the ITGC and retest cost of the one AI application the design uses. Under Assumption~\ref{ass:acct-itgc} that cost is carried by every application in $\mathcal U^*_{\mathrm{AI}}$, with its retest included in \eqref{eq:acct-itgc}, so the sum over $\mathcal U^*_{\mathrm{AI}}$ replaces it; with one application the two coincide. The other terms of \eqref{eq:acct-amin} are unchanged, and $C^{\mathrm{tml}}$ from (d) is added because it is a recurring cost of the reporting obligation that no other term holds. Layers in one application share its identities, permissions, changes, retests, operations and integrations, hence one line; separate applications have separate ones. An application in $\mathcal U^*_{\mathrm{AI}}$ through (S), (I) or (P) alone carries its line by Assumption~\ref{ass:acct-itgc}, whatever the control design. Lines are costs, so they do not enter $P_o$ and leave the admissible set of Definition~\ref{def:acct-admissible} unchanged. With one application per AI layer the cost of a design is $K(d)+g(n_{\mathrm{AI}})$ with $g(n)=(n-1)^+C^{\mathrm{itgc}}_g$ nondecreasing; if $d$ minimizes $K$ and $d'$ minimizes $K+g$ over the same set, $K(d')+g(n'_{\mathrm{AI}})\le K(d)+g(n_{\mathrm{AI}})$ and $K(d)\le K(d')$ give $g(n'_{\mathrm{AI}})\le g(n_{\mathrm{AI}})$, so $n'_{\mathrm{AI}}\le\max\{n_{\mathrm{AI}},1\}$, and the same argument with $g(n)=n\,\Delta C$ shows that a higher line never raises the minimizing number of AI layers (up to ties). [6c, 6d, 6e]

(c) The change in annual control cost on the success path is the sum of the lines added less those retired, plus the change in $C^{\mathrm{tml}}$, by Assumption~\ref{ass:acct-itgc} and (b). If upgrade dates form a stationary point process with intensity $1/T_{\mathrm{up}}$ independent of the cutover date (for example, dates spaced by $T_{\mathrm{up}}$ with a phase uniform on $[0,T_{\mathrm{up}})$), Campbell's theorem gives the expected present value of the retests on $[T_{\mathrm{par}},T_h]$ as $\int_{T_{\mathrm{par}}}^{T_h}e^{-rt}C^{\mathrm{val}}_g/T_{\mathrm{up}}\,dt$, the value of the flow. If the clock starts at cutover, the retests fall at $T_{\mathrm{par}}+kT_{\mathrm{up}}$ for $k=1,\dots,n_{\mathrm{up}}$, and the geometric sum $\sum_{k=1}^{n_{\mathrm{up}}}C^{\mathrm{val}}_ge^{-r(T_{\mathrm{par}}+kT_{\mathrm{up}})}$ is \eqref{eq:acct-pvup}. [6f, 6g]

(d) Incidents start in close windows at rate $\nu_{\mathrm{cw}}=n_{\mathrm{cl}}\nu_{\mathrm{out}}T_{\mathrm{cw}}$ a year (a thinned Poisson process), each with an independent duration $D\sim\mathrm{Exp}$ of mean $\bar D$. Without a fallback, an incident costs $c_{\mathrm{del}}D+L_{\mathrm{late}}\mathbf 1\{D>T_{\mathrm{sl}}\}$, with expectation $c_{\mathrm{del}}\bar D+L_{\mathrm{late}}e^{-T_{\mathrm{sl}}/\bar D}$. With one, the close is delayed by $\min\{D,t_{\mathrm{sw}}\}\le T_{\mathrm{sl}}$, so no deadline is missed, and the fallback runs for $(D-t_{\mathrm{sw}})^+$; since $\mathbb E\min\{D,t\}=\bar D(1-e^{-t/\bar D})$ and $\mathbb E(D-t)^+=\bar De^{-t/\bar D}$, the expected cost is as stated, plus $C^{\mathrm{fb}}$. Then $C^{\mathrm{tml}}_{\mathrm{wait}}-C^{\mathrm{tml}}_{\mathrm{fb}}=\nu_{\mathrm{cw}}\Delta^{\mathrm{fb}}-C^{\mathrm{fb}}$, which is positive if and only if $\Delta^{\mathrm{fb}}>0$ and $\nu_{\mathrm{out}}>\nu^*$. The firm chooses the cheaper, so $C^{\mathrm{tml}}=\min\{C^{\mathrm{tml}}_{\mathrm{wait}},C^{\mathrm{tml}}_{\mathrm{fb}}\}$. [6h, 6i, 6j] \qed
